\documentclass[10pt,journal,compsoc]{IEEEtran}

% *** CITATION PACKAGES ***
\ifCLASSOPTIONcompsoc
  \usepackage[nocompress]{cite}
\else
  \usepackage{cite}
\fi

% *** GRAPHICS ***
\ifCLASSINFOpdf
  \usepackage[pdftex]{graphicx}
  \graphicspath{{./figuras/}}
  \DeclareGraphicsExtensions{.pdf,.jpeg,.png}
\else
  \usepackage[dvips]{graphicx}
\fi

% *** MATH ***
\usepackage{amsmath}
\interdisplaylinepenalty=2500

\hyphenation{op-tical net-works semi-conduc-tor}

\begin{document}

\title{SiGAD: Sistema Interactivo de Guidance Contextual Explicativo\\
Alineado al Ciclo PPDAC para Fortalecer la Alfabetización\\
de Datos en Usuarios Novatos}

\author{Javier~Wilber~Quispe~Rojas
\IEEEcompsocitemizethanks{\IEEEcompsocthanksitem J. W. Quispe Rojas está con la
Escuela Profesional de Ciencia de la Computación, Universidad Nacional de San
Agustín de Arequipa (UNSA), Perú.
\IEEEcompsocthanksitem Asesora: Ana María Cuadros Valdivia, Escuela Profesional
de Ciencia de la Computación, UNSA.}%
\thanks{Manuscrito en preparación, 2026.}}

\markboth{Journal of \LaTeX\ Class Files,~Vol.~14, No.~8, August~2015}%
{Quispe Rojas \MakeLowercase{\textit{et al.}}: SiGAD: Sistema Interactivo de Guidance Contextual}

\IEEEtitleabstractindextext{%
\begin{abstract}
Los sistemas de \textit{guidance} actuales en analítica visual y ciencia de
datos concentran su apoyo en fases puntuales del análisis —un procedimiento
estadístico o la lectura de una visualización— sin acompañar al usuario desde
la formulación de la pregunta hasta la interpretación final. Como
consecuencia, los usuarios novatos no solo cometen errores de sintaxis, sino
que fallan al elegir operaciones, manipular datos o interpretar resultados,
limitando el desarrollo de su alfabetización de datos. Este artículo presenta
SiGAD, un sistema interactivo de \textit{guidance} contextual explicativo
alineado al ciclo PPDAC (Problema, Plan, Datos, Análisis y Conclusión) que
combina cálculo determinista sobre los datos con un modelo de guiado
pedagógico adaptativo —organizado en los niveles orientar, dirigir y
prescribir— cuyo nivel de apoyo se gradúa y retira según el dominio estimado
del estudiante. Se describe la arquitectura funcional del sistema y se
propone su evaluación frente a un artículo de referencia y frente a un grupo
de control, en términos de la calidad percibida del guiado y del efecto
sobre el desempeño, la comprensión metodológica y la autonomía de usuarios
novatos.
\end{abstract}

\begin{IEEEkeywords}
Guidance contextual, alfabetización de datos, ciclo PPDAC, guiado pedagógico
adaptativo, andamiaje, analítica visual.
\end{IEEEkeywords}}

\maketitle
\IEEEdisplaynontitleabstractindextext
\IEEEpeerreviewmaketitle

\IEEEraisesectionheading{\section{Introducción}\label{sec:introduccion}}

\IEEEPARstart{L}{os} entornos interactivos han comenzado a ofrecer apoyo
directo al usuario durante el análisis de datos. GuidedStats recibe un
\textit{dataset} ya provisto por el usuario y lo guía a través de nueve pasos
centrados en verificar supuestos y ajustar un procedimiento estadístico,
exportando los resultados a código al finalizar \cite{zhang2024}. VIZTA, por
su parte, se define como un asistente que ayuda a los lectores a comprender
visualizaciones ya construidas \cite{wang2025}. En ambos casos, el apoyo se
concentra en una fase puntual del análisis —el procedimiento estadístico o la
lectura del gráfico— sin acompañar al usuario desde la formulación de la
pregunta hasta la interpretación final.

Estos avances responden a la necesidad de fortalecer la alfabetización de
datos, entendida como la capacidad de recolectar, gestionar, evaluar y
aplicar datos de manera crítica para formular y responder preguntas del
mundo real mediante un ciclo iterativo de obtención, preparación, análisis,
interpretación y comunicación \cite{wolff2016}. Este ciclo retoma el modelo
PPDAC (Problema, Plan, Datos, Análisis y Conclusión), propuesto originalmente
por Wild y Pfannkuch como marco del pensamiento estadístico
\cite{wildpfannkuch1999} y adaptado por Wolff et al. para estructurar la
alfabetización de datos \cite{wolff2016}. Revisiones críticas recientes
confirman que sigue siendo un área activa de discusión, sin una definición
única consolidada \cite{lund2025}.

Pese a esto, la alfabetización de datos se promueve como competencia
transversal desde hace más de una década, y aun así los estudiantes
universitarios que completan cursos introductorios de análisis siguen siendo
incapaces de vincular coherentemente sus decisiones a lo largo del proceso,
particularmente al relacionar la preparación de datos con la validez de sus
conclusiones \cite{wolff2016}. La evidencia reciente confirma que estas
dificultades no se limitan a errores de sintaxis: los novatos también
cometen errores conceptuales y metodológicos al seleccionar operaciones,
manipular datos o interpretar salidas, sin comprender plenamente las
consecuencias de una transformación sobre el resto del análisis
\cite{singh2024,neuwinger2025}.

Los marcos de \textit{guidance} en analítica visual, como el de Ceneda et
al., caracterizan el grado de asistencia que un sistema puede ofrecer —desde
orientar hasta prescribir una acción— \cite{ceneda2016}, pero se fundamentan
en teoría de interacción visual, sin apoyarse en principios de aprendizaje
que expliquen por qué o cuándo esa asistencia debería graduarse para
fortalecer la competencia del usuario y no solo resolver su tarea inmediata;
es decir, ofrecen un mecanismo técnico de asistencia sin un fundamento
pedagógico que lo sostenga. Estos avances muestran que el \textit{guidance}
puede incorporar explicaciones dentro del flujo de análisis, pero, como se
señaló, se concentran en técnicas o fases puntuales y carecen de dicho
fundamento.

Esta investigación estudia el \textit{guidance} como un mecanismo
pedagógicamente fundamentado que acompaña el ciclo completo de análisis,
haciendo explícitas las razones metodológicas y las consecuencias de las
decisiones entre fases. Los \textit{guidance} existentes, al enfocarse en
técnicas específicas \cite{zhang2024,wang2025}, no acompañan todo el proceso
de análisis, por lo que los novatos no solo cometen errores de sintaxis, sino
que fallan al elegir operaciones, manipular datos o interpretar resultados,
limitando su alfabetización de datos \cite{wolff2016,singh2024,neuwinger2025}.

Para cerrar esta brecha se propone SiGAD (Sistema Interactivo de
\textit{Guidance} contextual explicativo Alineado al ciclo PPDAC para
fortalecer la Alfabetización de Datos), un sistema que combina cálculo
determinista sobre los datos con un modelo de guiado pedagógico
adaptativo —organizado en los niveles orientar, dirigir y prescribir— y con
generación de lenguaje natural para explicar, en cada fase del ciclo PPDAC,
por qué una decisión metodológica es adecuada o riesgosa. A diferencia de los
marcos de \textit{guidance} puramente técnicos, el nivel de apoyo en SiGAD se
ajusta según el dominio estimado del estudiante, replicando el principio de
retiro progresivo del andamiaje (\textit{scaffolding}) propuesto por Wood,
Bruner y Ross \cite{woodbrunerross1976} dentro de la Zona de Desarrollo
Próximo de Vygotsky \cite{vygotsky1978}.

Las contribuciones de este trabajo son las siguientes:
\begin{itemize}
\item Se diseña la arquitectura funcional de un sistema de guiado interactivo
que cubre las fases Problema, Plan, Datos, Análisis y Conclusión del ciclo
PPDAC.
\item Se implementa un modelo de guiado pedagógico adaptativo que integra el
intento autónomo inicial y los niveles orientar, dirigir y prescribir, con
reglas de escalamiento y retirada de apoyo basadas en el dominio estimado del
estudiante.
\item Se compara el sistema propuesto con el método de un artículo de
referencia, aplicando al grupo experimental una encuesta basada en los
criterios de adaptabilidad, oportunidad, explicabilidad y pertinencia del
apoyo.
\item Se evalúa el efecto del sistema en la alfabetización en ciencia de
datos, comparando el desempeño, la comprensión metodológica y la autonomía
entre estudiantes que usan el sistema y estudiantes que no.
\end{itemize}

Esta investigación se orienta, en consecuencia, por las siguientes preguntas:
¿cómo diseñar la arquitectura funcional de un sistema de guiado interactivo
que cubra las cinco fases del ciclo PPDAC (PI1)?; ¿cómo implementar un modelo
de guiado pedagógico adaptativo que gradúe y retire el apoyo (orientar,
dirigir, prescribir) según el dominio estimado del estudiante (PI2)?; ¿cómo
perciben los usuarios novatos la calidad del guiado proporcionado por el
sistema propuesto en términos de adaptabilidad, oportunidad, explicabilidad y
pertinencia del apoyo (PI3)?; y ¿qué diferencias se observan en el
desempeño, la comprensión metodológica y la autonomía de los usuarios
novatos que utilizaron el sistema propuesto frente a los que no lo utilizaron
(PI4)?

El resto de este artículo se organiza de la siguiente manera. La Sección~II
revisa los trabajos relacionados y posiciona a SiGAD frente a sistemas
educativos de alfabetización de datos, sistemas de alfabetización visual y
sistemas de \textit{guidance} en analítica visual. La Sección~III describe la
arquitectura y el modelo de guiado adaptativo propuestos. La Sección~IV
presenta el diseño experimental y los resultados de la evaluación.
Finalmente, la Sección~V discute las conclusiones y el trabajo futuro.
% NOTA: las Secciones II-V se agregan en un paso posterior; por ahora este
% archivo solo contiene Título, Resumen e Introducción.

\section*{Agradecimientos}
Los autores agradecen a la Escuela Profesional de Ciencia de la Computación
de la Universidad Nacional de San Agustín.

\begin{thebibliography}{99}

\bibitem{wolff2016}
A. Wolff, D. Gooch, J. J. Cavero~Montaner, U. Rashid, and G. Kortuem,
``Creating an Understanding of Data Literacy for a Data-driven Society,''
\emph{The Journal of Community Informatics}, vol.~12, no.~3, 2016.

\bibitem{wildpfannkuch1999}
C.~J. Wild and M. Pfannkuch, ``Statistical Thinking in Empirical Enquiry,''
\emph{International Statistical Review}, vol.~67, no.~3, pp. 223--248, 1999.

\bibitem{lund2025}
[Lund et al., 2025 — completar referencia con datos exactos del artículo].

\bibitem{singh2024}
[Singh et al., 2024 — completar referencia].

\bibitem{neuwinger2025}
[Neuwinger and Riehle, 2025 — completar referencia].

\bibitem{ceneda2016}
D. Ceneda, T. Gschwandtner, T. May, S. Miksch, H.-J. Schulz, M. Streit, and
C. Tominski, ``Characterizing Guidance in Visual Analytics,'' \emph{IEEE
Transactions on Visualization and Computer Graphics}, vol.~23, no.~1,
pp. 111--120, 2016.

\bibitem{zhang2024}
[Zhang et al., 2024 — GuidedStats — completar referencia].

\bibitem{wang2025}
[Wang et al., 2025 — VIZTA — completar referencia].

\bibitem{woodbrunerross1976}
D. Wood, J.~S. Bruner, and G. Ross, ``The Role of Tutoring in Problem
Solving,'' \emph{Journal of Child Psychology and Psychiatry}, vol.~17,
no.~2, pp. 89--100, 1976.

\bibitem{vygotsky1978}
L.~S. Vygotsky, \emph{Mind in Society: The Development of Higher
Psychological Processes}. Cambridge, MA: Harvard University Press, 1978.

\end{thebibliography}

\end{document}
